% =============================================================================
\section{Introduction}\label{sec:intro}

Put a \emph{Drosophila} larva on a dish of plain agar, with nothing to smell, see or eat, and it will not stay put. It crawls in straight runs separated by turns: many short runs and, every now and then, a very long one, with lengths spread over two decades in a way that resembles a power law~\cite{Sims2019}. Termite workers in a Petri dish~\cite{Castiblanco2019}, T cells in living tissue~\cite{Jerison2020} and marine predators in the open ocean~\cite{Humphries2010} leave records of the same kind. A walk made of straight flights whose lengths $\ell$ follow $p(\ell)\propto\ell^{-\mu}$, with $1<\mu\le3$, is called a L\'evy walk~\cite{Viswanathan1999,Zaburdaev2015}, and whether real organisms perform L\'evy walks, and why, has been argued about for two decades~\cite{Edwards2007,Pyke2015,Reynolds2018}.

Most of that argument concerns the histogram of flight lengths. This paper is about their order. Does a long flight make the next one longer? Does a walker remember its previous flights, and if so, how many? And if it does, is the exponent read from the histogram still the exponent of the rule that the walker follows at each turn? The textbook L\'evy walk has no memory: it is a renewal process, in which every flight is drawn afresh from the same law. A power law, however, only counts long flights. It does not say whether they come in runs, and runs of long flights change how far and how fast a walker spreads~\cite{Aghion2021}. When walkers move in groups a second question joins the first: when a walker speeds up or turns, is it following its own past or the past of the others?

Questions of this kind are inverse problems, and maximum entropy (MaxEnt) is a natural way to pose them~\cite{Jaynes1957}. Among all the probability laws that reproduce a chosen set of measured averages, MaxEnt selects the one that assumes nothing else. When the averages are those of some functions $f_k$, the selected law is an exponential family, $p\propto\exp(\sum_k\lambda_kf_k)$, formally a Boltzmann weight whose Lagrange multipliers $\lambda_k$ act as fields and couplings; read in this way, the correlations measured in flocks of birds and in populations of neurons have been turned into networks of effective interactions~\cite{Bialek2012,Mora2016,Roudi2011,Nguyen2017}. Maximizing the entropy of whole trajectories rather than of single configurations, the principle of maximum caliber~\cite{Presse2013,Ge2012}, carries the idea over to dynamics. Constraints on correlations at different times have been used to predict the spiking of neurons~\cite{Marre2009} and the motion of birds in a flock~\cite{Cavagna2014} and, in the dynamic-entropy formalism of Agliari \emph{et al.}~\cite{Agliari2024}, to infer delayed interactions between migrating cells from the directions in which they move.

We apply this program to the lengths of flights. Its static half is old: fixing only the mean of $\ln\ell$, MaxEnt returns a power law~\cite{Montroll1983}. The dynamic half is not. Fixing also the correlations between $\ln\ell_t$ and $\ln\ell_{t-\tau}$, for lags up to a depth $K$, turns the sequence of flights into a chain of spins coupled in time, and the depth that the data select tells a renewal walker ($K=0$) from a Markovian one ($K=1$) and from one with a longer memory. The chain is a small statistical-mechanics problem of its own. With one lag it can be solved exactly. Exchanging long and short flights, $\ell\to\lmin\lmax/\ell$, maps the chain onto itself with different parameters, and on the line of parameters that this duality leaves unchanged the flights condense at the largest available scale through a first-order transition, with the logarithmic range of the flights in the role of system size. Short of the transition, memory makes the histogram heavier-tailed than the rule that generates it. In a group, the same construction separates the correlation of a walker with the past into a part due to its own past and a part due to the past of the others, and the collective speed acquires a critical point.

On open data, memory turns out to be common. Larvae crawling without sensory cues remember one to four flights back, and they remember most when their brain lobes and sensory neurons have been ablated. Termites remember about four flights back, and the wall of their dish, which they follow closely, does not explain it. In groups, wild baboons align with where their nearest neighbors, those ahead above all, were heading 10 to 30\,s earlier, and homing pigeons are coupled to the whole of their small flocks. Zebrafish T cells, which we use as a negative control, do not interact at all, although an analysis of their group average would say otherwise.

We do not ask whether these walks are optimal searches, whether a power law describes their flights better than a log-normal law (over one or two decades it often does not, Appendix~\ref{app:data}), or which circuits in the nervous system produce the memory. The question is narrower: given a trajectory, is there memory, how deep is it, and whose is it?

Section~\ref{sec:model} builds the model, for one walker and for many, and Sec.~\ref{sec:statmech} works out its physics. Section~\ref{sec:inference} turns the model into an estimator and Sec.~\ref{sec:data} applies it to larvae, termites, baboons, pigeons and T cells. Section~\ref{sec:discussion} discusses what the results mean for L\'evy walks. Proofs, validations on synthetic data and the controls on the real ones are collected in the Appendices.
